\section{Complete finite recipe for the five coefficients}
\label{app:coefficient-recipe}

This appendix specifies all five fixed integers in the main polynomial.
Its powers, finite sums, and finite case distinctions concern only fixed
coefficients: none contains a Diophantine unknown. The complete window
list and the binary expansions of the resulting integers are not
materialized. The correction bit defined below is likewise retained as
an exact finite arithmetic expression, rather than claimed to have been
numerically evaluated.

The literal K215 table and its normalized transition function are fixed
in Section~\ref{sec:controller}. For this appendix, denote that normalized
function by
\[
 \delta(q,s)=(q',s',D),\qquad
 q,q'\in\{0,\ldots,434\},\quad
 s,s'\in\{0,\ldots,9\},\quad D\in\{-1,0,1\}.
\]
Its start state is $0$, its erase state is $1$, its absorbing halt is $4$,
and its blank symbol is $0$. In particular $\delta(0,1)=(1,2,0)$.
The archive supplies the entire 4,350-row table, so this use of $\delta$
does not introduce an unspecified universal-machine index.

\subsection{The complete ordered tile alphabet and local predicate}
\label{app:tile-predicate}

Write $\bot$ for an absent payload, phase, or head, as appropriate.
A tile has four fields $(v,h,\pi,\varphi)$: two boundary bits, a payload,
and a phase. A payload is $(s,q)$, with a tape symbol $s$ and a state
$q$ or $\bot$. The ordered tile alphabet is as follows:
\begin{enumerate}
\item Tile $0$ is $\mathsf V=(1,0,\bot,\bot)$ and tile $1$ is
$\mathsf H=(0,1,\bot,\bot)$.
\item For $s=0,\ldots,9$ and $q=\bot,0,\ldots,434$, in this order,
include $(0,0,(s,q),\bot)$. Its tile index is
\[
 2+436s+\eta(q),\qquad
 \eta(\bot)=0,\quad \eta(q)=q+1\quad(q\ne\bot).
\]
\item Tiles $4362,4363,4364,4365$ have phases
$\mathsf L,\mathsf I,\mathsf Q,\mathsf R$, respectively. Their boundary
bits are zero, and their respective payloads are
\[
 (0,\bot),\quad(1,\bot),\quad(1,0),\quad(0,\bot).
\]
\end{enumerate}
Thus the alphabet has $a_c=4366$ elements. All the fields of every tile
are fixed finite data.

Here is the complete cross predicate used to select windows. Let
$t,t_{\ell},t_r,t_d,t_u$ be the center, left, right, down, and up tiles.
Subscripts on fields mean that they are read from the corresponding tile.
First impose
\begin{equation}\label{app:boundary-clause}
 v=v_u,\qquad vv_r=0,\qquad hh_u=0,\qquad
 (v=v_r=0\ \Longrightarrow\ h=h_r).
\end{equation}
If the center is $\mathsf V$ or $\mathsf H$, these are its entire tests.
For a nonboundary center, impose all the following additional conditions.
\begin{enumerate}
\item If $h_d=1$, its phase belongs to
$\{\mathsf L,\mathsf I,\mathsf Q,\mathsf R\}$ and its payload is the
one prescribed above. A left $\mathsf V$ requires phase $\mathsf L$;
a right $\mathsf V$ requires phase $\mathsf R$. If the right neighbor is
not $\mathsf V$, the two phases must form one of the six pairs
\[
 (\mathsf L,\mathsf L),\ (\mathsf L,\mathsf I),\
 (\mathsf I,\mathsf I),\ (\mathsf I,\mathsf Q),\
 (\mathsf Q,\mathsf R),\ (\mathsf R,\mathsf R).
\]
If $h_d=0$, the center phase must instead be $\bot$.
\item If the center carries a head, its transition may not move left
through a left $\mathsf V$ or right through a right $\mathsf V$.
If $h_u=1$, that head, if present, must be in halt state $4$.
\item If $h_u=0$, each non-$\mathsf V$ left or right neighbor must have
a payload. Treat a neighboring $\mathsf V$ as payload $(0,\bot)$.
The next center symbol is its current symbol when headless, and its written
symbol under $\delta$ otherwise. Form the list of incoming states from
the left head moving right, the center head staying, and the right head
moving left. There must be at most one such state. The up payload must
be the next center symbol together with that state, or with $\bot$ if
the list is empty.
\end{enumerate}
This is the literal predicate in the standalone
\path{artifacts/compiler/finite_window_constants.py}; it retains every
initialization, collision, no-escape, and final-row test of the inherited
helical relation~\cite{CompilerDetails}.

A $3\times3$ window is allowed exactly when its central cross satisfies
these tests. Its four corner entries are unrestricted. Write entries in
row-major order, with row $0$ supplying the down neighbor and row $2$
supplying the up neighbor. The two distinguished windows, expressed
entirely in tile indices, are
\[
 \mathsf S=
 \begin{pmatrix}
 1&1&1\\4363&4363&4364\\438&438&876
 \end{pmatrix},\qquad
 \mathsf E=
 \begin{pmatrix}
 1&1&1\\4362&4363&4363\\2&438&438
 \end{pmatrix}.
\]
Both pass the predicate; they are the Start and End windows. The ordered
list $w_0,\ldots,w_{k_c-1}$ begins with $\mathsf S,\mathsf E$ and then
contains every other allowed nine-tuple in lexicographic tile-index order.
This defines a terminating enumeration with at most $4366^9$ predicate
tests, independent of whether the machine halts. Its exact size is
\begin{equation}\label{app:window-count}
 k_c=113175624520853754856480928432.
\end{equation}
The exact compressed census is the sum of
$311471256024587$ allowed central crosses, each with $4366^4$ choices
of corners. It is a finite partition count, not a sample estimate.

\subsection{Evaluated fixed layout parameters}

The following parameters determine the native positions, clause bands,
and radices. The subscript $c$ distinguishes compiler parameters from the
Pell variables in the main polynomial.
\begin{center}
\small
\begin{tabular}{@{}lr@{}}
\toprule
Parameter & Fixed value\\
\midrule
$a_c$ & $4366$\\
$\ell_c$ & $97$\\
$p_c$ & $56587812260426877428240418274$\\
$m_c$ & $113175624520853754856480967730$\\
$U_c$ & $113175624520853754856480993923$\\
$H_c$ & $5771956850563541497680530703172$\\
$T_{c,1}$ & $11883440574689644259930504379381$\\
$T_{c,2}$ & $17994924298815747022180478051224$\\
$g_c$ & $21050666160878798403305464887146$\\
\bottomrule
\end{tabular}
\end{center}
Put
\begin{equation}\label{app:fixed-radices}
 A_c=2^{\ell_c},\qquad b=5^{44},\qquad L_c=5^{45},\qquad
 d=bL_c=5^{89},\qquad R_0=2^b,\qquad B=R_0^{L_c}=2^d.
\end{equation}
These values reproduce the deterministic layout rule:
\begin{align*}
 \ell_c&=\max\{2,\operatorname{bitlength}(k_c+1)\},&
 m_0&=2\ell_c+30a_c+8,\\
 p_c&=\max\left\{0,\left\lceil
       \frac{k_c+9a_c+4-m_0}{2}\right\rceil\right\},&
 m_c&=m_0+2p_c,\\
 U_c&=m_c+6a_c-3,&H_c&=51U_c+3a_c+1,\\
 T_{c,1}&=105U_c+4a_c+2,&T_{c,2}&=159U_c+4a_c+3,\\
 g_c&=186U_c+4a_c+4.
\end{align*}
The least power of five at least $213U_c+4a_c+5$ is $L_c$.
The radix choice is justified after the clause coefficients are defined.

\subsection{Native positions and all clause coefficients}

For $r,s\in\{0,1,2\}$ and $u\in\{0,\ldots,a_c-1\}$, define
\begin{align}
 j_c(r,s,u)&=1+3a_cr+a_cs+u,\label{app:clause-index}\\
 e_c(r,s,u)&=k_c+3a_c+(2-r)3a_c+(2-s)a_c+u.
 \label{app:payload-position}
\end{align}
There is exactly one ignored dummy position,
\[
 e_* = k_c+12a_c =113175624520853754856480980824.
\]
The complete finite native position set is
\begin{align}
 \mathcal E_c={}&\{0,\ldots,k_c-1\}
 \cup\{e_c(r,s,u):0\leq r,s\leq2,\ 0\leq u<a_c\}
 \nonumber\\
 &\cup\{e_*\}\cup\{U_c,3U_c,9U_c,27U_c\}.
 \label{app:native-positions}
\end{align}
Its members are distinct, and $|\mathcal E_c|=m_c+1$.

Define a nonnegative fixed coefficient $c_e$ at each position. For a
selector position $0\leq i<k_c$, set
\begin{equation}\label{app:selector-coefficient}
 c_i=1+
 \sum_{r=0}^{2}\sum_{s=0}^{2}
 A_c^{\,j_c(r,s,(w_i)_{3r+s})}
 +\sum_{h=1}^{4}A_c^{\,9a_c+h}.
\end{equation}
Here $(w_i)_t$ is the tile index at entry $t$ of the already specified
fixed window $w_i$. At the remaining positions put
\begin{align}
 c_{e_c(r,s,u)}&=A_c^{j_c(r,s,u)},&c_{e_*}&=0,\nonumber\\
 c_{U_c}&=A_c^{9a_c+1},&c_{3U_c}&=A_c^{9a_c+2},\nonumber\\
 c_{9U_c}&=A_c^{9a_c+3},&c_{27U_c}&=A_c^{9a_c+4}.
 \label{app:other-coefficients}
\end{align}
The clause mask, including its fixed zero-expression padding clauses, is
\begin{equation}\label{app:clause-mask}
 \mu_c=A_c-2+\sum_{j=1}^{9a_c+4+p_c}A_c^j.
\end{equation}
Omitting the padding summands would change the mask population and would
not be the inherited compiler instance.

For clarity, the least-power-of-five radix rule can now be stated exactly.
Let $C_\Sigma=\sum_{e\in\mathcal E_c}c_e$. Its target is
\[
 T_{\mathrm{radix}}=
 \max\{4(m_c+3)(2C_\Sigma+6)+8,\ 2\mu_c+4,\ 16\}.
\]
Choose $b$ to be the least power of five at least
$\operatorname{bitlength}(T_{\mathrm{radix}}-1)$.
For this instance,
\[
 C_\Sigma<(14k_c+9a_c+5)A_c^{9a_c+4}<2^{3812007},
\]
so the first member of the maximum is below $2^{3812108}$.
The second member dominates. Its exact bit-length contribution is
\begin{align*}
 \operatorname{bitlength}(2\mu_c+3)
 &=\ell_c(9a_c+4+p_c)+2\\
 &=5489017789261407110539324384486.
\end{align*}
Its least power-of-five majorant is $5^{44}$, as asserted
in~\eqref{app:fixed-radices}. No giant power has to be expanded to verify
these inequalities.

\subsection{Convolution coefficients and the fixed correction bit}

The uncorrected center coefficient and the right coefficient are
\begin{align}
 D_{C,0}={}&R_0^{3a_c}+R_0^{H_c+a_c}
             +R_0^{8U_c}+R_0^{24U_c}\nonumber\\
 &+\sum_{e\in\mathcal E_c}c_e
       \bigl(R_0^{T_{c,1}-e}+R_0^{T_{c,2}-e}\bigr),
 \label{app:center-convolution}\\
 D_R={}&R_0^{H_c}.
 \label{app:right-convolution}
\end{align}
All exponents are nonnegative fixed integers. Define the single fixed bit
\begin{equation}\label{app:epsilon}
 \epsilon_c=
 \begin{cases}
 1,&2D_{C,0}-D_R\equiv0\pmod5,\\
 0,&2D_{C,0}-D_R\not\equiv0\pmod5,
 \end{cases}
 \qquad
 D_C=D_{C,0}+\epsilon_c R_0^{g_c}.
\end{equation}
This is a finite arithmetic decision on uniquely defined fixed integers.
It is not a Diophantine variable, an unspecified choice of coefficient,
or a query about consistency. Since $b$ is a power of five,
$R_0\equiv2\pmod5$. The added monomial has nonzero residue after
multiplication by $2$, so the definition guarantees
\[
 2D_C-D_R\not\equiv0\pmod5.
\]
It can in principle be evaluated using modular exponentiation and the
fixed finite sums. This packet does not claim to have evaluated
$\epsilon_c$ or the window-dependent sum in~\eqref{app:center-convolution}.

\subsection{Native masks, shifted field mask, and final ports}

Define the old field mask by the exact finite expression
\begin{align}
 M_{F,0}={}&\mu_c\bigl(R_0^{T_{c,1}}+R_0^{T_{c,2}}\bigr)
 \nonumber\\
 &+\sum_{r=0}^{1}\sum_{s=0}^{2}\sum_{u=0}^{a_c-1}
       R_0^{e_c(r,s,u)}
 \nonumber\\
 &+\sum_{r=0}^{2}\sum_{s=0}^{1}\sum_{u=0}^{a_c-1}
       R_0^{H_c+e_c(r,s,u)}
       +R_0^{9U_c}+R_0^{27U_c}.
 \label{app:old-field-mask}
\end{align}
The modified half-binomial masks are
\begin{align}
 M_C&=B-1-\sum_{\substack{e\in\mathcal E_c\\e\ne1}}R_0^e
           -2R_0^{e_*},\label{app:native-mask}\\
 M_F&=M_{F,0}+4.\label{app:field-mask}
\end{align}
The first modification frees the second bit at the ignored dummy position;
the second adds the low verification bit. These are the modifications of
the reviewed compiler export~\cite{CompilerDetails}. In particular
\begin{gather*}
 0<M_C,M_F<B-1,\qquad
 M_C\equiv2\pmod4,\qquad M_F\equiv4\pmod8,\\
 \pop(M_C)+\pop(M_F)=d.
\end{gather*}

The five coefficients supplied to the main polynomial are now exactly
\begin{equation}\label{app:five-final-coefficients}
 \boxed{
 \begin{aligned}
 \beta&=B-1,&K&=D_C+BD_R,\\
 M&=M_C,&N&=M_F+B-1,\\
 W&=R_0B^2=2^{\,2\cdot5^{89}+5^{44}}.
 \end{aligned}}
\end{equation}
The source coefficient $N$ is already shifted. It is \emph{not} the
native one-cell mask $M_F$; replacing one by the other changes the
polynomial. The endpoint $W$ inserts the End selector at cyclic cell $2$,
giving the three-I initialization run and four unary ones already proved
equivalent to blank K215.

\subsection{Explicitness and the representation boundary}

Every summation bound, finite set, field lookup, and exponent in this
appendix is fixed before the sixteen unknowns of the main polynomial
are chosen. Substitution of~\eqref{app:five-final-coefficients} therefore
produces one ordinary integer-coefficient polynomial. No variable power,
bounded universal quantifier, sequence parameter, or acceptance test
remains in that polynomial.

The compressed representation is essential for practical delivery:
$B$ alone has $5^{89}+1$ binary digits. This coefficient height is not
included in the 70-operation evaluation count. The count assumes access
to the five fixed numerals, as stated in the main theorem; it does not
assert that their construction or storage is inexpensive. The companion
\path{artifacts/compiler/COEFFICIENTS.md}, source code, and pinned
manifests give the same recipe in machine-readable and prose forms.

All compiler formulas in this appendix are the lazy form of the retained
construction
\[
 \texttt{new\_constants(compile\_windows(windows,a))}.
\]
The new explicit controller instantiation does not replace the inherited
proof of that compiler, and the exact finite specification must not be
confused with a claim that all its giant integers have been expanded.
